% Answers to Chapter 2, translated from sv/back/facit/02.tex.

\facitavsnitt{c2}{Chapter 2}{Algebra}
\begin{facit}

\svar{1.1} $R_{\text{TOT}} = 1\,\Omega$

\svar{1.2} \sv{a} $460\,\text{W}$ \sv{b} $40\,\text{W}$

\svar{1.3} $U_{\text{out}} = 4\,\text{V}$

\svar{2.1} \sv{a} $3I_1 + 10I_2$ \sv{b} $3R_1 + 2R_2$ \sv{c} $7I^2 + 3I - 1$
\sv{d} Both expressions give $33$.

\svar{2.2} \sv{a} $U_R = RI_1 + RI_2 + RI_3$; every term has the unit $\Omega \cdot \text{A} =
\text{V}$. \sv{b} $10\,\text{V}$ \sv{c} $2 + 3 + 5 = 10\,\text{V}$, the same answer.

\svar{2.3} \sv{a} $3I(2 + 3R)$ \sv{b} $(P_1 + P_2)(P_1 - P_2)$ \sv{c} $(U - 5)^2$
\sv{d} $(2R + 1)(2R - 1)$

\svar{2.4} \sv{a} $P_{\text{TOT}} = (R_1 + R_2 + R_3)I^2$ \sv{b} $1.25\,\text{W}$
\sv{c} $P_1 = 0.30\,\text{W}$, $P_2 = 0.45\,\text{W}$, $P_3 = 0.50\,\text{W}$, sum
$1.25\,\text{W}$. The factorised form took the least work: one multiplication by $I^2$ instead of
three.

\svar{2.5} \sv{a} Take out $\frac{1}{R}$; what is left is $U_1^2 - U_2^2$, a difference of two
squares, which factorises by that rule. \sv{b} $22\,\text{W}$
\sv{c} $P_1 = 72\,\text{W}$ and $P_2 = 50\,\text{W}$; the difference is $22\,\text{W}$.

\svar{2.6} \sv{a} $P = RI_0^2 + 2RI_0\,\Delta I + R\,\Delta I^2$ \sv{b} $44.1\,\text{W}$
\sv{c} $\Delta P = 4.1\,\text{W}$ \sv{d} $0.1\,\text{W}$, about $2.4\,\%$ of $\Delta P$.

\svar{2.7} \sv{a} $RI^2 = RI \cdot I$, and $RI = U$ is the voltage across the resistor by Ohm's
law, so $P = UI$. \sv{b} $R \cdot \frac{U^2}{R^2} = \frac{U^2}{R}$
\sv{c} $I = 2\,\text{A}$; all three forms give $24\,\text{W}$.

\svar{3.1} \sv{a} Three \sv{b} $7$ and $-4$ respectively \sv{c} $9$
\sv{d} $7$ and $R^2$ (or $7$, $R$ and $R$) \sv{e} $29$

\svar{3.2} \sv{a} $3I + 14$ \sv{b} $U + 7$ \sv{c} $4R_1 + 3$ \sv{d} $5$

\svar{3.3} \sv{a} $R^2 + 8R + 15$ \sv{b} $2I^2 + 7I - 4$ \sv{c} $U^2 - 4$
\sv{d} $9R^2 - 12R + 4$ \sv{e} $7I$

\svar{3.4} \sv{a} $I^2 + 8I + 16$ \sv{b} $R^2 - 12R + 36$ \sv{c} $4U^2 + 12U + 9$
\sv{d} $25 - 10I + I^2$ \sv{e} $(3 + 4)^2 = 49$ but $3^2 + 4^2 = 25$; the difference is twice the
product, $2ab = 24$.

\svar{3.5} \sv{a} $9996$ \sv{b} $1575$ \sv{c} $3.99$ \sv{d} $200$ \sv{e} $80\,\text{W}$

\svar{3.6} \sv{a} $4(2R + 3)$ \sv{b} $RI(I + 1)$ \sv{c} $(U + 7)(U - 7)$ \sv{d} $(I + 6)^2$
\sv{e} $(3R + 4)(3R - 4)$ \sv{f} $2(U + 2)(U - 2)$

\svar{3.7} \sv{a} $2I$ \sv{b} $R - 3$ \sv{c} $U + 2$ \sv{d} $I - 5$ \sv{e} $2R + 1$

\svar{3.8} \sv{a} $12\,\text{V}$ \sv{b} $18\,\text{W}$ \sv{c} $18\,\text{W}$, the same answer.
\sv{d} The power quadruples, to $72\,\text{W}$: $R(2I)^2 = 4RI^2$.

\svar{3.9} \sv{a} The fractions have the same denominator, and the numerators add up to the
denominator $R_1 + R_2$, so $U_1 + U_2 = U_{\text{in}}$.
\sv{b} In the quotient, $U_{\text{in}}$ and the denominator $R_1 + R_2$ cancel, leaving
$\frac{R_1}{R_2}$. \sv{c} $U_1 = \frac{U_{\text{in}}}{2}$
\sv{d} $U_1 = 3\,\text{V}$, $U_2 = 12\,\text{V}$

\svar{3.10} \sv{a} Twice the product is missing: $R^2 + 4R + 4$.
\sv{b} The factor must multiply both terms: $3I - 6$.
\sv{c} A fraction can only be cancelled by factors: $\frac{R}{4} + 1$.
\sv{d} The minus sign changes the sign of every term: $-U + 3$.
\sv{e} The variables are multiplied too: $6R^2$.
\sv{f} The exponent applies to the whole bracket: $4I^2$.

\svar[Test yourself]{T} \sv{1} $(3R + 2)(3R - 2)$ \sv{2} $5I + 7$
\sv{3} $20\,\text{W}$

\end{facit}
